\clearpage
\phantomsection% 
\addcontentsline{toc}{chapter}{RINGKASAN}
\thispagestyle{fancy}

\begin{center}
    \large \bfseries \MakeUppercase{Ringkasan}\\
    \normalsize \normalfont {\thetitle}\\
    \normalsize \normalfont {\theauthor}\\
    \bigskip
\end{center}

\normalsize \normalfont \justifying
Prevalensi \textit{stunting} di Indonesia merupakan masalah kesehatan masyarakat yang dipengaruhi oleh banyak faktor. Hubungan antarindikator kesehatan ini sering kali bersifat non-linear, sehingga pemodelan linear konvensional kurang sesuai untuk menganalisis akar masalahnya. Penelitian ini bertujuan untuk membangun model estimasi prevalensi \textit{stunting} tingkat provinsi di Indonesia, mengidentifikasi faktor determinan utama, serta menganalisis titik ambang batas pada setiap indikator menggunakan data Survei Status Gizi Indonesia (SSGI) tahun 2024.

Penelitian ini menggunakan pendekatan \textit{Machine Learning} dengan algoritma \textit{Random Forest Regression}. Parameter model dioptimalkan melalui metode \textit{Grid Search} dan dievaluasi menggunakan \textit{Leave-One-Out Cross-Validation} (LOOCV). Pemrosesan data melibatkan 54 fitur determinan yang berpotensi memengaruhi angka \textit{stunting}. Selanjutnya, untuk menginterpretasikan hasil prediksi model \textit{Random Forest} yang bersifat \textit{black-box}, diterapkan metode \textit{SHapley Additive exPlanations} (SHAP) guna menganalisis kontribusi setiap fitur terhadap variabel target.

Hasil pengujian menunjukkan bahwa algoritma \textit{Random Forest Regression} dapat memodelkan data agregat kesehatan dengan baik. Evaluasi model menghasilkan nilai koefisien determinasi ($R^2$) sebesar 0,9090, \textit{Mean Absolute Error} (MAE) sebesar 1,40, \textit{Root Mean Squared Error} (RMSE) sebesar 1,77, dan \textit{Mean Absolute Percentage Error} (MAPE) sebesar 7,22\%. Hasil ini menunjukkan bahwa model yang dibangun cukup layak digunakan untuk mengestimasi angka prevalensi \textit{stunting} tingkat provinsi.

Berdasarkan visualisasi \textit{SHAP Feature Importance} dan \textit{Summary Plot}, faktor determinan utama yang memengaruhi prevalensi \textit{stunting} antarprovinsi berkaitan dengan kelengkapan asupan gizi balita, kondisi kehamilan, dan fasilitas sanitasi. Fitur yang paling berdampak meliputi ketiadaan pemenuhan \textit{Minimum Dietary Diversity} (MDD) dan \textit{Minimum Milk Feeding Frequency} (MMFF), persentase usia ibu berisiko saat hamil, ketiadaan sanitasi layak, serta rendahnya cakupan Pemberian Makanan Tambahan (PMT).

Lebih lanjut, analisis \textit{SHAP Dependence Plot} memperlihatkan adanya titik ambang batas risiko pada beberapa indikator utama. Beberapa titik batas yang ditemukan antara lain ketiadaan MDD di angka 55\%, ketiadaan MMFF di 22\%, usia ibu berisiko sebesar 14\%, dan fasilitas sanitasi berisiko sebesar 10\%. Apabila persentase indikator di suatu provinsi melewati nilai tersebut, risiko kenaikan prevalensi \textit{stunting} akan cenderung meningkat.

Secara keseluruhan, pemodelan \textit{Random Forest Regression} yang digabungkan dengan interpretasi SHAP dapat digunakan sebagai alat bantu analisis faktor penyebab \textit{stunting}. Nilai ambang batas yang ditemukan dapat dijadikan informasi tambahan bagi instansi terkait dalam mengevaluasi program perbaikan gizi dan sanitasi lingkungan. Untuk penelitian selanjutnya, disarankan menggunakan data dengan cakupan tingkat kabupaten/kota serta melakukan perbandingan dengan metode interpretasi lain seperti LIME untuk memverifikasi konsistensi penjelasan model.

\vfill
\clearpage